\documentclass[12pt]{article}
\usepackage[utf8]{inputenc}
\usepackage[T1]{fontenc}
\usepackage{amsmath,amssymb}
\usepackage{geometry}
\geometry{margin=2.5cm}

\begin{document}

\noindent Ele sunt adevărate prop:

\begin{enumerate}
\item[1.] $\big(a,(b,c)\big) = \big((a,b),c\big)$
\item[2.] $(ca,cb) = c(a,b)$.
\item[3.] dacă $(a,b)=1$ şi $(a,c)=1 \Rightarrow (a,bc)=1$.

\noindent generalizare

\[
(a,b_i)=1,\ i=\overline{1,n} \ \Rightarrow \ (a,\,b_1b_2\cdots b_n)=1.
\]

\item[4.] $a\mid bc$, $(a,b)=1 \ \Rightarrow \ a\mid c$

\item[5.] $\left.\begin{aligned} a&\mid c\\ b&\mid c\\ (a,b)&=1\end{aligned}\right\} \Rightarrow ab\mid c$

\noindent generalizare: \quad $m_i\mid c$, \ $(m_i,m_j)=1 \ \Rightarrow \ m_1m_2\cdots m_n\mid c$

\item[6.] $(a,b)=1 \Rightarrow$ [continuare neclară în original: fragmentul vizibil e ,,$c,ab$''.]\footnote{proprietatea 6 e incompletă/neclară în original -- posibil să implice $(ab,c)$, dar formularea exactă nu e sigură.}
\end{enumerate}

\bigskip
\noindent (1) evident.

\noindent $(b,c)=d$.\footnote{rolul exact al acestei notaţii în demonstraţie e neclar în original.}

\bigskip
\noindent (3). \ $1=(a,b)=(a,c)$

\bigskip
\noindent (4) \ $a\mid bc \Rightarrow (a,bc)=a$

\[
(a,c) \overset{?}{=} \big((a,bc),c\big) = \big(a,(bc,c)\big)
\]
\[
\Rightarrow (a,c)=a \ \Rightarrow \ a\mid c
\]
\noindent\footnote{pasul algebric din mijloc e foarte compactat în original şi nu poate fi reconstituit cu certitudine deplină -- te rog verifică-l.}

\bigskip
\noindent (5). \ $(ab,c)=$ \ [pagina se întrerupe aici.]

\bigskip
\noindent Fie $p$ = nr.\ nat.\ ireductibil $\Rightarrow p$=prim

\noindent\underline{Dem}: R.A.\ (reducere la absurd) \ $p\neq$prim $\Rightarrow \exists\,a,b\in\mathbb{N}$ a.î.

\[
p\mid ab,\quad p\nmid a,\quad p\nmid b
\]
\[
\Rightarrow (p,a)=1 \text{ şi } (p,b)=1 \overset{(3)}{\Longrightarrow} (p,ab)=1
\]
\[
\Rightarrow p=1 \quad \times.
\]

\end{document}
